\section*{Open Questions (Q1--Q5)}

\begin{description}
\item[Q1. DSB complexity distribution vs total number as driver of LET-scaling.]
\textit{Basis:} Friedland 2012 encodes complexity implicitly via the slow-channel
saturation term but never publishes the DSB-complexity distribution shape
parameters that PARTRAC generates; the paper's phenomenology cannot be
decomposed into a yield vs complexity contribution.
\textit{Next steps:} (a) run Geant4-DNA at UICGPU for 5 ion species
(H/He/C/N/Fe) $\times$ 5 energies each; extract per-DSB complexity histograms;
(b) fit a two-component NHEJ Gillespie model where the slow rate is a function
of DSB complexity bin, not aggregate LET; (c) cross-validate on KURBUC and
RITRACKS on identical target geometries.

\item[Q2. Cross-code comparison of Friedland's NHEJ rules on PARTRAC vs Geant4-DNA vs KURBUC.]
\textit{Basis:} PARTRAC is closed; no independent group has run a rigorous
cross-code test of Friedland's exact NHEJ formulation. Published paper-to-paper
comparisons conflate physics differences with chemistry/biology differences.
\textit{Next steps:} (a) implement the Friedland 2012 NHEJ rules
(fast/slow rejoining + complexity-dependent slow capacity + labile-site delay)
as an open Gillespie library; (b) drive it with Geant4-DNA-generated then
KURBUC-generated DSB end lists; diff rejoining kinetics at 15~min / 1~h / 4~h
/ 24~h; (c) publish the library with a test-suite.

\item[Q3. Chromatin geometry sensitivity of the two-rate sufficiency claim.]
\textit{Basis:} Friedland 2012 assumes a stylised chromatin model inside
PARTRAC. Hi-C-derived polymer models show DSB proximity statistics depend
strongly on chromatin state; part of what Friedland attributes to biochemistry
may be geometry.
\textit{Next steps:} (a) use a Hi-C-informed polymer simulator (chromatin3D or
OpenMiChroM) to sample DSB proximity graphs for euchromatin- and
heterochromatin-enriched regions; (b) drive the open Gillespie library from Q2
with those proximity graphs and refit fast/slow rate constants; (c) test
whether the two-component sufficiency claim breaks (needs a third component or
a state-dependent rate).

\item[Q4. Sub-DSB damage spectrum contribution to the ``labile-site'' population.]
\textit{Basis:} Friedland 2012 introduces a labile-site amplitude but does not
decompose it by underlying lesion type. Distinguishing enzyme kinetics from
lesion-conversion kinetics is a key open question for both photon and ion RBE
modelling.
\textit{Next steps:} (a) run Geant4-DNA with chemistry stage on to enumerate
sub-DSB clustered lesion densities per LET; (b) fit a three-population NHEJ
model (prompt-DSB / labile-converted DSB / clustered-lesion converted DSB) to
photon and ion kinetics; (c) test whether the fitted labile amplitude scales
with predicted clustered-lesion density across the LET grid.

\item[Q5. Extension to the FLASH ultra-high-dose-rate regime.]
\textit{Basis:} FLASH ($>$40 Gy/s) may alter DSB production and repair via
transient oxygen depletion or radical recombination, phenomena absent from
PARTRAC's conventional kinetics. The paper predates FLASH and does not
comment.
\textit{Next steps:} (a) literature-scan FLASH-vs-conventional DSB rejoining
measurements (Vozenin, Favaudon, Loo groups); (b) add a dose-rate-dependent
oxygen-depletion and radical-recombination term to the open Gillespie library
and test whether the two-component + labile skeleton still fits FLASH
kinetics; (c) cross-check against Geant4-DNA-Chem and transient $p$O$_2$
measurements.
\end{description}
